\documentclass[conference]{IEEEtran}
\usepackage{graphicx}
\usepackage{amsmath}
\usepackage{float}
\usepackage{siunitx}
\usepackage{url}
\usepackage{hyperref} % This makes the links clickable
\begin{document}

% --- Cover Page Elements ---
\title{Design and Performance Analysis of a Grid-Connected Solar PV System\\[0.5cm]
\large 0406320 Solar PV Systems --- Group Project (SAM)}
\author{
    \IEEEauthorblockN{Oumar Ibrahim}
    \IEEEauthorblockA{College of Computing and Informatics, University of Sharjah\\
    Sharjah, United Arab Emirates\\
    Email: \href{mailto:u22200741@sharjah.ac.ae}{u22200741@sharjah.ac.ae}}
}

\maketitle

\begin{abstract}
This report details the design and simulation of a roof-mounted photovoltaic (PV) system for the University of Sharjah (UOS) male sports complex. Utilizing the System Advisor Model (SAM), the system is configured across two tilted roof subarrays to maximize available footprint. The document covers site assessment, hardware selection, and a comprehensive loss analysis. The simulation results demonstrate that the proposed design achieves an optimized balance between annual energy yield and system efficiency, providing a viable renewable energy solution for the facility.\end{abstract}

\begin{IEEEkeywords}
Photovoltaic systems, System Advisor Model (SAM), Solar energy, Inverters, Grid-connected PV.
\end{IEEEkeywords}

\section{Site Assessment}
The proposed PV system is located at the UOS male sports complex in Sharjah (Coordinates: 25.284503799178378, 55.478224746237565). The design utilizes the two tilted sides of the facility's roof, with PV modules installed flush to the roof at an assumed tilt of $5^{\circ}$ for each side \cite{google_earth}. 

% Google Earth Screenshot
\begin{figure}[htbp]
    \centering
    \includegraphics[width=0.45\textwidth]{googleearthazimuth2.png}
    \caption{Google Earth screenshot showing roof measurements. The measured roof azimuth angles are 36.22$^{\circ}$ and 216.22$^{\circ}$. The measured total roof area of 2,228.98 m$^2$.}
    \label{fig:google_earth}
\end{figure}

\section{Initial Calculations}
Prior to utilizing SAM \cite{sam_software}, initial manual calculations were performed to determine the maximum system capacity, appropriate inverter sizing, and string configurations based on Standard Test Conditions (STC) of 1000 W/m$^2$. The selected hardware for this commercial-scale design includes Trina Solar TSM-330 modules (17\% efficiency, $V_{oc}$ = 46.2V) and SMA Sunny Tripower 60kW string inverters (Min $V_{DC}$ = 415V, Max $V_{DC}$ = 980V).

\subsection{System Size Capacity}
Using the measured total roof area of 2,228.98 m$^2$ (approximately 1,114.49 m$^2$ available per subarray), the maximum system capacity is estimated:
\begin{equation}
    C_{\text{array}} = A_{\text{roof}} \times \eta_{\text{module}} \times 1000 \text{ W/m}^2
\end{equation}
\begin{equation}
    C_{\text{array}} = 2228.98 \text{ m}^2 \times 0.17 \times 1000 \text{ W/m}^2 = 378,926 \text{ W}
\end{equation}
The maximum estimated array capacity is approximately 378.9 kW DC.
Dividing the estimated array capacity by the module maximum power yields the number of modules in our design:
\begin{equation}
    TotalModules = \frac{C_{\text{array}}}{P_{MP}} = \frac{378.9 kW}{329.6 W} \approx 1150 
\end{equation}
\subsection{Inverter Sizing}
Assuming a target commercial DC/AC ratio of approximately 1.2, the required inverter AC capacity is calculated as follows:
\begin{equation}
    \text{Inverter AC Capacity} = \frac{\text{System Size (DC W)}}{1.2}
\end{equation}
\begin{equation}
    \text{Inverter AC Capacity} = \frac{378,926 \text{ W}}{1.2} = 315,771 \text{ W}
\end{equation}
To optimize the physical array layout and perfectly fit the available roof space without overlapping the edges, the final design utilizes five SMA Sunny Tripower 60kW inverters, providing a total AC capacity of 300 kW. This yields a highly optimized design ratio.

\subsection{Module and String Sizing}
Based on the chosen module's open-circuit voltage ($V_{oc}$) of 46.2V and the selected SMA inverter's specific input voltage limits, the string constraints are determined:
\begin{equation}
    \text{Max modules/string} = \frac{\text{Inverter Max V}_{DC}}{\text{Module } V_{oc}} = \frac{980\text{V}}{46.2\text{V}} = 21.2
\end{equation}
\begin{equation}
    \text{Min modules/string} = \frac{\text{Inverter Min V}_{DC}}{\text{Module } V_{oc}} = \frac{415\text{V}}{46.2\text{V}} = 9.0
\end{equation}
To maximize the utilization of the available roof footprint, a configuration of 18 modules per string and 31 strings per subarray was selected. This arrangement results in a total of $1,116$ modules, providing the highest feasible energy density while remaining strictly within the physical capacity limit of $1,150$ modules

\section{Final System Design in SAM}
The calculated parameters were implemented in SAM. The design features two distinct subarrays to account for the two opposite-facing roof sides.

% System Design Screenshot
\begin{figure}[htbp]
    \centering
    \includegraphics[width=0.45\textwidth]{"System design.png"}
    \caption{SAM ``System Design'' page detailing the ratio of module per string, number of inverters, and the number of strings in parallel in a 2-subarray configuration.}
    \label{fig:system_design}
\end{figure}

% Result Summary Screenshot
\begin{figure}[htbp]
    \centering
    \includegraphics[width=0.45\textwidth]{results.png}
    \caption{SAM Result Summary page displaying the overall 629,127 kWh annual energy production and metrics.}
    \label{fig:result_summary}
\end{figure}

\section{System Conclusions and Performance}
Based on the SAM simulation using the downloaded Sharjah weather data, the system yields the following performance metrics:

\begin{itemize}
    \item \textbf{Total PV DC Power:} As shown in \autoref{fig:system_design}, The total installed capacity across all modules is 367.85 kW.
    \item \textbf{Peak System Power:} The maximum AC power generated by the system over the year is capped at 300 kW due to the inverter capacity limit, occurring consistently during the peak spring and summer months.
\end{itemize}


\begin{figure}[htbp]
    \centering
    \includegraphics[width=0.25\textwidth]{subarray.png}
    \caption{Gross DC power generated by each subarray.}
    \label{fig:subarray}
\end{figure}
\subsection{Production Comparisons}
\begin{itemize}
    \item \textbf{Subarray Productivity:}  \autoref{fig:subarray} illustrates the gross DC generated by each subarray where the Southwest-facing roof (Subarray 2 at 216.22$^{\circ}$) is more productive. It generates 346,997 Gross DC kWh annually, whereas the Northeast-facing roof (Subarray 1 at 36.22$^{\circ}$) generates 327,751 Gross DC kWh.
    \item \textbf{Monthly Extremes:} As shown in \autoref{fig:monthly}, the most productive month is May with approximately 66,000 kWh generated. Conversely, the least productive month is January, yielding approximately 39,000 kWh. \autoref{fig:time_series} illustrates peak generation from March to June, driven by high UAE spring irradiance. During this period, the system consistently operates near its 300~kW AC limit before extreme summer temperatures begin to degrade module efficiency.
\end{itemize}

% Figure 1 Screenshot
\begin{figure}[htbp]
    \centering
    \includegraphics[width=0.25\textwidth]{monthly.png}
    \caption{AC energy production across the months of the year.}
    \label{fig:monthly}
\end{figure}
% Time Series Screenshot
\begin{figure}[htbp]
    \centering
    \includegraphics[width=0.45\textwidth]{timeseries.png}
    \caption{Time series graph from SAM showing system generated power (kW) across the year.}
    \label{fig:time_series}
\end{figure}
\section{Loss Analysis and Optimization}

\subsection{System Sizing Macro Results}
The System Sizing macro was executed to validate the array-to-inverter pairing and check for clipping or sizing issues.

% Table 2 Screenshot
\begin{figure}[htbp]
    \centering
    \includegraphics[width=0.45\textwidth]{Macro.png}
    \caption{Table 2 results from the System Sizing macro showing array and inverter matching.}
    \label{fig:macro_table2}
\end{figure}

% Figure 1 Screenshot
\begin{figure}[htbp]
    \centering
    \includegraphics[width=0.45\textwidth]{FIgure1.png}
    \caption{Figure 1 results from the System Sizing macro illustrating power distributions.}
    \label{fig:macro_figure1}
\end{figure}

Based on these macro results illustrated by \autoref{fig:macro_table2}, \autoref{fig:macro_figure1}    and the main simulation data in \autoref{fig:result_summary}, the design does not exhibit significant sizing losses. The inverter power clipping loss is a negligible 0.002\%, indicating that the chosen 1.23 DC/AC ratio using five 60kW inverters perfectly handles the array's peak output without unnecessarily bottlenecking the generated energy.

% Losses Page Screenshot
\begin{figure}[htbp]
    \centering
    \includegraphics[width=0.25\textwidth]{losses.png}
    \caption{SAM ``Losses'' page overview detailing the percentage breakdown of system inefficiencies.}
    \label{fig:sam_losses}
\end{figure}


\subsection{System Losses Analysis}


Reviewing the detailed losses page shown in \autoref{fig:sam_losses}, the primary performance reductions include Module deviation from STC (7.999\%) and Soiling (5.0\%). The high STC deviation is heavily expected due to the extreme ambient temperatures in the UAE, which degrade the efficiency of silicon cells. To improve the overall performance, the following strategies are suggested:

\begin{enumerate}
    \item Implement a bi-weekly or monthly automated cleaning schedule to mitigate the high 5\% soiling losses caused by dust and sand accumulation in the Sharjah environment.
    \item Select a premium PV module with a lower temperature coefficient, such as N-type silicon or advanced Heterojunction (HJT) cells. Recent industry analysis indicates that N-type and HJT modules maintain significantly higher efficiency at temperatures exceeding 25$^{\circ}$C compared to standard P-type modules, effectively reducing heat-related power degradation in desert climates. [3],[4].
\end{enumerate}

\section{Conclusion}
This project successfully designed a commercial-scale 367.85 kW roof-mounted PV system for the UOS male sports complex. By carefully analyzing the available roof geometry and employing a highly optimized 1.23 DC/AC inverter ratio, the dual-subarray design maximizes the facility's footprint to generate 629,127 kWh annually. While regional factors like high heat and soiling introduce notable losses, the system operates with negligible (0.002\%) inverter clipping, making it a highly efficient and viable renewable energy solution for the university.


\begin{thebibliography}{00}

\bibitem{sam_software}
National Renewable Energy Laboratory (NREL), \emph{System Advisor Model (SAM)}, Version 2025.4.16. [Online]. Available: \url{https://sam.nrel.gov/}


\bibitem{google_earth}
Google, \emph{Google Earth Pro}, Version 7.3. [Online]. Available: \url{https://www.google.com/earth/}

\bibitem{maysunsolar}
Maysun Solar, \emph{PV Module Performance Comparison in High Temperature}. [Online]. Available: \url{https://www.maysunsolar.com/blog-pv-module-performance-comparison-in-high-temperature/}

\bibitem{astadvancedenergy}
AST Advanced Energy, \emph{TOPCon vs HJT vs Back Contact Solar PV Modules}. [Online]. Available: \url{https://astadvancedenergy.com/blog/hello-world/}

\end{thebibliography}


\end{document}
